\documentclass[conference]{IEEEtran}
\usepackage{cite}
\usepackage{amsmath,amssymb,amsfonts}
\usepackage{booktabs}
\usepackage{array}
\usepackage{multirow}
\usepackage{xcolor}
\usepackage{url}
\usepackage{microtype}
\usepackage{balance}
\setlength{\tabcolsep}{3pt}
\newcolumntype{L}[1]{>{\raggedright\arraybackslash}p{#1}}
\newcommand{\specscope}{\textsc{Speciescope}}
\newcommand{\designonly}{\textsc{Design-Only}}

\begin{document}

\title{Counting Life Without Conflating Units: \\
\specscope{} for Auditable Estimates of Earth Species Richness}

\author{\IEEEauthorblockN{Four-Stage Research Workflow}
\IEEEauthorblockA{\textit{Survey - Idea - Experiment Design - Author}\\
\textnormal{\designonly{} research artifact}\\
Earth biodiversity estimation}}

\maketitle

\begin{abstract}
``How many species are there on Earth?'' is often presented as a disagreement among headline numbers. It is primarily an estimand problem: named taxa, eukaryotic taxonomic species, lineage hypotheses, barcode clusters, operational taxonomic units, and amplicon sequence variants do not necessarily count the same biological entities. This design-only paper synthesizes the key evidence boundary and introduces \specscope{}: Species-Concept, Evidence, Coverage, Integration, Estimation, and Scope Evaluation. The framework requires an estimand passport before any global number is reported, models detection and unseen richness within defined domains, separates molecular discovery from taxonomic assertion, tests extrapolation on held-out strata, and blocks arithmetic aggregation when units are incompatible. A five-stage protocol, preregistered endpoints, data governance rules, and failure-oriented validation plan are specified. The central conclusion is not a replacement total. It is that an estimate becomes scientifically interpretable only when its species concept, spatial and environmental frame, evidence rule, estimator, uncertainty, and integration status are jointly reported.
\end{abstract}

\begin{IEEEkeywords}
biodiversity estimation, species delimitation, sample coverage, environmental DNA, microbial diversity, taxonomic data integration
\end{IEEEkeywords}

\section{Introduction}

The number of species on Earth is a foundational quantity for ecology, evolution, conservation, and biological discovery. It also resists reduction to a single settled value. May emphasized the large uncertainty inherited from incomplete observation and uneven taxonomic knowledge decades ago~\cite{may1988}. Modern sequencing, digital registries, and global observation networks have increased the rate at which diversity can be detected, yet they also introduced several measurement units whose relationship to a species is conditional rather than automatic.

Two influential estimates illustrate the issue. Mora \emph{et al.} inferred approximately 8.7 million global eukaryotic species, with a reported standard error of 1.3 million, by extrapolating a pattern in higher taxonomic classification~\cite{mora2011}. This is a transparent model-based estimate for a specified eukaryotic target, not a direct census of all biological entities. Locey and Lennon combined global community data, a dominance scaling relationship, and a lognormal model to predict up to $10^{12}$ microbial ``species''~\cite{locey2016}. The latter result foregrounds enormous microbial diversity, but it has a distinct measurement and extrapolation boundary. Treating these values as rival readings of one common instrument, or adding them, conflates target populations and units.

This paper asks a more tractable scientific question: \emph{given a declared species concept, taxonomic scope, geographic and environmental sampling frame, time window, evidence standard, and estimator, what richness interval and uncertainty are supported?} It proposes \specscope{} as a reproducible answer architecture. The objective is not to manufacture a more dramatic global number. The objective is to make every legitimate number auditable, comparable where appropriate, and explicitly non-additive where necessary.

The contribution is fourfold. First, the paper formalizes the estimand passport that every richness claim should carry. Second, it separates three inferential layers---detection, delimitation, and extrapolation---whose errors are often hidden in a single species-count figure. Third, it specifies a five-stage study protocol that uses specimen data, registries, environmental DNA (eDNA), and statistical estimators without assuming that they produce interchangeable units. Fourth, it defines integration gates that permit a combined total only after overlap, unit mapping, and uncertainty representation have passed a predeclared audit.

This is a \designonly{} research report. It contains no new collecting, sequencing, taxonomic revision, or numerical Earth-richness result. All quantities introduced below define planned measurements, estimators, validation metrics, or literature boundaries rather than empirical outcomes of this study.

\section{Why the Headline Range Is Not One Numerical Dispute}

\subsection{Species concepts and evidence}

The word ``species'' has multiple operational definitions because speciation is a process and organisms vary in reproduction, geography, asexuality, hybridization, morphology, and genomic history. de Queiroz frames many practical concepts as evidential criteria for separately evolving metapopulation lineages~\cite{dequeiroz2007}. This does not eliminate disagreement, but it provides a disciplined interpretation: morphology, reproductive isolation, ecological differentiation, and phylogenetic evidence are alternative evidence streams, not necessarily interchangeable definitions.

Cryptic diversity shows why no single evidence type suffices. Morphologically similar lineages may be evolutionarily distinct, while shallow genetic splits may reflect population structure rather than species boundaries~\cite{bickford2007,leliaert2014}. Molecular methods therefore create a productive tension. They can reveal candidate diversity at a scale impossible for traditional cataloguing alone, but a sequence cluster becomes a taxonomic claim only under a stated delimitation protocol. Single-locus approaches such as generalized mixed Yule-coalescent procedures may be useful provisional tools, while their accuracy depends on divergence, population structure, and model assumptions~\cite{fujisawa2013}.

\begin{table*}[t]
\caption{Common biodiversity counting units and their appropriate interpretation.}
\label{tab:units}
\centering
\footnotesize
\begin{tabular}{L{0.16\textwidth} L{0.27\textwidth} L{0.25\textwidth} L{0.24\textwidth}}
\toprule
\textbf{Unit} & \textbf{Scientific meaning} & \textbf{Strength} & \textbf{Non-negotiable caveat}\\
\midrule
Accepted named taxon & A nomenclatural concept in a versioned taxonomic registry & Traceable names, synonym history, and expert revision & Counts depend on description effort and the selected registry release.\\
Specimen-supported lineage hypothesis & Candidate separately evolving lineage supported by declared evidence & Can combine morphology, ecology, geography, and genomics & Evidence sufficiency must be explicit and vouchers must be retained.\\
Barcode Index Number (BIN)-like unit & Persistent sequence-based registry unit for a marker locus & Reproducible discovery and record linkage & A barcode unit is not automatically a biological or formally named species~\cite{ratnasingham2013}.\\
OTU or ASV & Bioinformatic cluster or resolved sequence feature & High-throughput, auditable molecular detection & Richness changes with marker, denoising, clustering, and curation choices.\\
Coverage-adjusted richness & Statistical estimate for a declared sampling universe & Quantifies observed and unseen components & It is conditional on detection, sample frame, and rare-tail assumptions.\\
\bottomrule
\end{tabular}
\end{table*}

Table~\ref{tab:units} distinguishes units that are routinely compressed into ``species.'' The distinction matters most at the interface between well-curated macrobiota and microbial communities. High-throughput sequencing can inventory molecular units across many environments, but sequencing error, primers, extraction, reference gaps, and thresholding influence the observed table. Post-clustering curation can substantially affect biodiversity estimates~\cite{froslev2017}. Neither the productivity of molecular surveys nor the expertise of taxonomists licenses an undocumented conversion from one unit to another.

\subsection{What published global estimates establish}

The estimate by Mora \emph{et al.} is often cited as a total for ``life on Earth,'' although its stated principal estimate is approximately 8.7 million \emph{eukaryotic} species~\cite{mora2011}. Its higher-taxon extrapolation is valuable because it sets out assumptions and an uncertainty expression. It should be reported as an estimate for that target, together with the scope and period of the underlying taxonomy. It should not be interpreted as an exhaustive count of microbial lineages, viruses, or every potential operational unit.

The microbial result of Locey and Lennon is likewise informative within its model. Their compilation spanned approximately 35,000 sites and 5.6 million species entries, and their scaling-law/lognormal combination predicted up to one trillion microbial species~\cite{locey2016}. The claim's interpretation depends on the data integration, abundance distributions, global extrapolation, and microbial delineation adopted by the study. The rare biosphere and its environmental heterogeneity make the model scientifically important; they do not make all molecular clusters one-to-one equivalents of taxonomic species.

The public range from several million to one trillion therefore contains at least three statements: a description backlog, an eukaryotic model estimate, and a microbial diversity extrapolation. It is not evidence that one research team must be wrong. The correct response is to declare the estimand and then test the associated design and model.

\section{The \specscope{} Framework}

\subsection{Estimand passport}

The first stage, S0, prevents a calculation from being performed before its target is defined. Let a richness estimand be the tuple
\begin{equation}
\mathcal{E}=\left(U,C,D,G,T,P,\mathcal{A}\right),
\label{eq:estimand}
\end{equation}
where $U$ is the counting unit, $C$ the species concept or delimitation criterion, $D$ the taxonomic domain, $G$ the geographic and environmental frame, $T$ the time window, $P$ the observation and processing protocol, and $\mathcal{A}$ the estimator and analysis version. Equation~\eqref{eq:estimand} makes a numerical result incomplete if any component is absent. A database release, primer set, taxonomy, and bioinformatic pipeline are not implementation footnotes; they define what was counted.

The passport is recorded before inspection of a planned richness estimate. It lists inclusion and exclusion rules, taxonomic rank, registry version, sample-frame map, access restrictions, unit crosswalks, and the exact meaning of ``unseen.'' A named-taxon count may use $U=$ accepted registry taxon and a global nomenclatural scope. A microbial survey may instead use $U=$ denoised ASV within a specific marker, habitat, and processing protocol. Both are valuable, but their passports expose why arithmetic summation is normally invalid.

\begin{table}[t]
\caption{\specscope{} five-stage protocol and main fail-fast rule.}
\label{tab:stages}
\centering
\footnotesize
\begin{tabular}{L{0.14\columnwidth} L{0.36\columnwidth} L{0.40\columnwidth}}
\toprule
\textbf{Stage} & \textbf{Deliverable} & \textbf{Fail-fast rule}\\
\midrule
S0 Scope & Estimand passport and unit map & Stop if unit, domain, or time window is not declared.\\
S1 Evidence & Controlled detection and candidate-delimitation record & Stop promotion of a claim when controls fail or linkage evidence is absent.\\
S2 Coverage & Detection, occupancy, and coverage diagnostics & Stop extrapolation if the sample frame is unknown or coverage is inadequate.\\
S3 Estimation & Interval, sensitivity distribution, and held-out calibration & Reject a preferred model if it is unstable or miscalibrated.\\
S4 Integration & Versioned ledger and additivity decision & Report in parallel if units, overlaps, or uncertainty models are incompatible.\\
\bottomrule
\end{tabular}
\end{table}

Table~\ref{tab:stages} presents the S0--S4 workflow. It is intentionally designed to yield a scientifically useful non-result: the decision that two estimates must remain parallel. Such a decision prevents a false all-life total and identifies precisely what evidence would be required to revisit integration.

\begin{figure*}[t]
\centering
\fbox{\begin{minipage}{0.94\textwidth}
\centering
\vspace{0.45em}
\begin{tabular}{@{}c@{\hspace{0.9em}}c@{\hspace{0.9em}}c@{\hspace{0.9em}}c@{\hspace{0.9em}}c@{}}
\fbox{\parbox[c][2.3em][c]{0.15\textwidth}{\centering\footnotesize\textbf{S0 Scope}\\Estimand passport}} &
$\Longrightarrow$ &
\fbox{\parbox[c][2.3em][c]{0.15\textwidth}{\centering\footnotesize\textbf{S1 Evidence}\\Controls and delimitation}} &
$\Longrightarrow$ &
\fbox{\parbox[c][2.3em][c]{0.15\textwidth}{\centering\footnotesize\textbf{S2 Coverage}\\Detection and unseen units}}\\\noalign{\vskip1.25em}
\fbox{\parbox[c][2.3em][c]{0.15\textwidth}{\centering\footnotesize\textbf{S3 Estimation}\\Intervals and sensitivity}} &
$\Longrightarrow$ &
\fbox{\parbox[c][2.3em][c]{0.15\textwidth}{\centering\footnotesize\textbf{S4 Integration}\\Add only when gates pass}} &
$\Longrightarrow$ &
\fbox{\parbox[c][2.3em][c]{0.30\textwidth}{\centering\footnotesize\textbf{Versioned output}\\Scope-specific estimates or a parallel non-additive ledger}}
\end{tabular}
\vspace{0.65em}
\end{minipage}}
\caption{\specscope{} turns an apparent one-number question into a gated evidence workflow. Failure at a gate narrows or pauses the claim; it does not justify an unsupported aggregate. An editable draw.io source with the same logic accompanies the report.}
\label{fig:workflow}
\end{figure*}

Fig.~\ref{fig:workflow} shows the information flow implied by the protocol. Scope is upstream of every other decision: coverage calculations only have a scientific meaning after the target unit and frame are frozen. Evidence and coverage then constrain estimation, and integration is deliberately last. The figure's two possible outputs emphasize that a transparent ledger of non-additive estimates is often more informative than a single global sum.

\subsection{Evidence, detection, and delimitation}

S1 records how a candidate unit was observed and how it was delimited. For a repeated detection design, $Y_{ijks}$ denotes whether unit $i$ is detected in visit $j$, method $k$, and sample $s$. A minimal observation model is
\begin{equation}
Y_{ijks}\,|\,Z_{is}\sim\mathrm{Bernoulli}\left(Z_{is}p_{ijks}\right),
\label{eq:detection}
\end{equation}
where $Z_{is}$ is latent occurrence and $p_{ijks}$ is method- and context-specific detection probability. Equation~\eqref{eq:detection} distinguishes non-detection from absence. It also creates a place to test laboratory batch, observer, season, substrate, sequencing depth, and primer effects instead of silently absorbing them into richness.

The latent occurrence probability can vary among environmental strata $h$ as
\begin{equation}
Z_{is}\sim\mathrm{Bernoulli}(\psi_{ih}), \qquad s\in h,
\label{eq:occupancy}
\end{equation}
where $\psi_{ih}$ is occupancy for candidate $i$ in stratum $h$. Equation~\eqref{eq:occupancy} supports replicated occupancy analyses, provided that the design has replication sufficient to estimate detection. It does not convert an ASV into a species. Rather, it estimates the observation process for the declared unit.

Specimen-based observations follow a parallel provenance rule: collection event, identifier, voucher repository, georeference precision, imagery where appropriate, nomenclatural concept, and revision history are recorded. Molecular observations require field, extraction, and PCR blanks; positive mock communities; technical replicates; raw read retention; denoising and filtering parameters; and reference-database versions. eDNA and metabarcoding can greatly increase detection of difficult-to-observe diversity~\cite{deiner2017,olds2016}, but controls are prerequisites for inferential promotion, not optional metadata.

\begin{table*}[t]
\caption{Detection and delimitation confounds with predeclared responses.}
\label{tab:confounds}
\centering
\footnotesize
\begin{tabular}{L{0.18\textwidth} L{0.25\textwidth} L{0.27\textwidth} L{0.22\textwidth}}
\toprule
\textbf{Confound} & \textbf{How it distorts a count} & \textbf{Measurement response} & \textbf{Decision recorded in passport}\\
\midrule
False negative & A present unit is missed, inflating the unseen component & Repeat visits, alternate methods, occupancy modelling & Detection model, visit structure, and covariates\\
Contamination or index leakage & A non-present molecular unit appears observed & Blanks, mock communities, batch tracking, independent replication & Threshold and exclusion rule\\
Primer/extraction bias & Some lineages are preferentially amplified or retained & Mock communities, multi-marker comparison, standard operating procedures & Marker, primer, extraction, and filtering versions\\
Synonymy and taxonomic revision & Several names refer to one concept, or one name is split & Frozen registry release, concept identifiers, revision ledger & Source registry and synonym policy\\
Cryptic diversity & Morphological records collapse lineages & Vouchers, targeted genomic/morphological follow-up, uncertainty labels & Evidence threshold for lineage status\\
Cluster threshold & OTU richness changes under a bioinformatic choice & ASV retention, alternative thresholds, curation and sensitivity tests & Algorithm, parameters, and reference database\\
\bottomrule
\end{tabular}
\end{table*}

\section{Coverage and Richness Estimation Within a Defined Domain}

\subsection{Coverage before extrapolation}

A raw observed count is not comparable across samples with unequal completeness. Let $f_1$ and $f_2$ be the numbers of units observed once and twice in an abundance sample of size $n$. A simple sample-coverage diagnostic can be expressed as
\begin{equation}
\widehat{C}=1-\frac{f_1}{n}\frac{(n-1)f_1}{(n-1)f_1+2f_2}.
\label{eq:coverage}
\end{equation}
Equation~\eqref{eq:coverage} is a planning diagnostic, not a universal truth. It communicates how much of the abundance mass is represented by detected units and motivates coverage-based rather than equal-effort comparisons~\cite{chao2012}. Incidence formulations are substituted when samples record occurrence rather than individual abundance.

For an abundance-based lower-bound estimator, the observed richness $S_{\mathrm{obs}}$ can be augmented by an estimated unseen component,
\begin{equation}
\widehat{S}=S_{\mathrm{obs}}+\widehat{S}_{0}, \qquad
\widehat{S}_{0}=\frac{n-1}{n}\frac{f_1^2}{2\max(f_2,1)}.
\label{eq:unseen}
\end{equation}
Equation~\eqref{eq:unseen} illustrates why singletons and doubletons matter: the rare tail carries information about undetected units. Its denominator protection is a practical expression of instability when $f_2$ is sparse, not a reason to claim false precision. The final study will select an estimator family that matches abundance, incidence, and spatial dependence; it will never present Eq.~\eqref{eq:unseen} as a one-size-fits-all global estimator.

The design uses an environmental and taxonomic stratum index $h=1,\ldots,H$. Stratum-level estimates are retained before any aggregate is considered. A weighted domain estimate has the schematic form
\begin{equation}
\widehat{S}_{D}=\sum_{h=1}^{H}w_h\widehat{S}_{h}, \qquad \sum_{h=1}^{H}w_h=1,
\label{eq:stratified}
\end{equation}
where $w_h$ encodes a declared target-frame weight rather than mere sampling convenience. Equation~\eqref{eq:stratified} forces researchers to state whether the target is area-weighted, habitat-weighted, host-weighted, or another well-defined estimand. Convenience samples can supplement coverage but cannot silently define global prevalence.

\subsection{Sampling design}

The proposed field architecture is a stratified, rotating panel spanning terrestrial, freshwater, marine, subterranean, atmospheric, host-associated, and engineered environments. Each broad realm is subdivided by biome, depth or elevation, substrate, latitude, climate, host association, and anthropogenic pressure. A probability-sampling core protects inference; museum, monitoring, and citizen-science observations are incorporated as auxiliary data with an explicit effort model. Rare or under-sampled strata can be oversampled while retaining inclusion probabilities and target weights.

For macrobiota, the record is a specimen or directly observed organism linked to a collection event and a voucher whenever feasible. For molecular surveys, the record is a physical environmental sample linked to extraction, library, sequencing run, controls, and a versioned bioinformatic workflow. The two streams are deliberately coupled at calibration sites: expert inventory, barcodes or genome references, and eDNA are collected in the same spatiotemporal window. This creates an evidence bridge rather than assuming that a database match resolves every species boundary.

\begin{table}[t]
\caption{Primary endpoints and predeclared interpretation.}
\label{tab:endpoints}
\centering
\footnotesize
\begin{tabular}{L{0.32\columnwidth} L{0.54\columnwidth}}
\toprule
\textbf{Endpoint} & \textbf{Interpretation}\\
\midrule
Coverage-adjusted richness interval & Richness for one passport-defined domain and unit, not a universal total.\\
Held-out detection calibration & Whether predicted detection and occurrence agree with new replicated samples.\\
Richness calibration & Whether withheld-stratum observations are consistent with predeclared interval predictions.\\
Sensitivity distribution & Change induced by plausible marker, threshold, taxonomy, weight, and rare-tail choices.\\
Uncertainty decomposition & Relative contribution of detection, delimitation, coverage, modelling, and integration uncertainty.\\
Integration-gate decision & Integrable, parallel-but-non-additive, or insufficiently specified.\\
\bottomrule
\end{tabular}
\end{table}

\subsection{Molecular units and taxonomic promotion}

Molecular surveys should publish both discovery units and taxonomic assertions. BIN-like systems demonstrate the utility of persistent sequence registries for accelerating biodiversity discovery~\cite{ratnasingham2013}. However, gene trees can disagree with species trees because of incomplete lineage sorting, introgression, hybridization, and other processes~\cite{leliaert2014}. A proposed promotion ladder is: (i) validated molecular unit, (ii) candidate lineage hypothesis, (iii) taxonomically supported species concept, and (iv) accepted named taxon in a frozen registry. Each step records new evidence; none is implied by merely changing a column heading.

The rare microbial biosphere needs an additional safeguard. A microbial ASV, a 97\% OTU, a genome-resolved population, and a formally named prokaryotic species can encode different biological and analytical choices. A global microbial study must report its specific unit, markers or genomes, abundance and occupancy criteria, error controls, and scale of extrapolation. This makes a trillion-scale prediction comparable as a model claim, even if not additively comparable with an estimate of eukaryotic taxonomic species.

\section{Estimation, Sensitivity, and Integration}

\subsection{Model competition and sensitivity}

S3 evaluates several defensible models rather than selecting the most impressive total. Let $\widehat{S}(a)$ be the estimate under analytical specification $a$, including a delimitation threshold, reference version, detection model, rare-tail model, and stratum weighting. Define a normalized local sensitivity measure as
\begin{equation}
\Delta(a,a')=\frac{\left|\widehat{S}(a)-\widehat{S}(a')\right|}{\max\left(\widehat{S}(a),\epsilon\right)},
\label{eq:sensitivity}
\end{equation}
where $a'$ is a plausible alternative specification and $\epsilon>0$ avoids a degenerate denominator. Equation~\eqref{eq:sensitivity} must be reported as a distribution over a predeclared set of alternatives, not cherry-picked for a favorable value. Large sensitivity is a finding about measurement dependence; it is not a defect to be concealed.

The primary analysis is preregistered for each domain, with a limited family of alternative models. Candidate models must meet coverage and control requirements, yield uncertainty intervals, and be evaluated on held-out sites, times, or strata. Nested resampling must preserve spatial and phylogenetic dependence where those structures are relevant. Exploratory model searches are separately labelled, and any later confirmatory test uses independent data or a locked validation split.

Overall uncertainty is not reducible to a sampling standard error. A transparent decomposition is
\begin{equation}
V_{\mathrm{total}}=V_{\mathrm{det}}+V_{\mathrm{delim}}+V_{\mathrm{cov}}+V_{\mathrm{model}}+V_{\mathrm{int}}+2\sum_{r<q}\mathrm{Cov}(r,q),
\label{eq:variance}
\end{equation}
where the variance components respectively represent detection, delimitation, coverage, model, and integration uncertainty. Equation~\eqref{eq:variance} explicitly retains covariance terms rather than treating all sources as independent. When a component cannot be estimated, the report must mark it as unquantified and explain its likely direction and consequence.

\subsection{Integration gates}

The fundamental output of S4 is a structured vector of domain estimates,
\begin{equation}
\mathbf{S}=\left(\widehat{S}_{1}^{U_1,C_1},\widehat{S}_{2}^{U_2,C_2},\ldots,\widehat{S}_{m}^{U_m,C_m}\right),
\label{eq:vector}
\end{equation}
where the superscripts identify unit and criterion. Equation~\eqref{eq:vector} is intentionally not replaced by $\sum_m \widehat{S}_m$. Summation is permitted only if all components share an eligible mapping to a common target unit, have a defined overlap rule, cover a declared compatible frame, and include an uncertainty model for the mapping and overlap. Otherwise, the vector is the scientifically correct report.

\begin{table*}[t]
\caption{Integration gates for deciding whether domain estimates can be combined.}
\label{tab:gates}
\centering
\footnotesize
\begin{tabular}{L{0.20\textwidth} L{0.31\textwidth} L{0.22\textwidth} L{0.19\textwidth}}
\toprule
\textbf{Gate} & \textbf{Required evidence} & \textbf{Failure consequence} & \textbf{Permitted output}\\
\midrule
Unit mapping & Auditable crosswalk from each unit to a common target, with ambiguity labels & Cannot determine whether one entity is counted once or many times & Parallel estimates only\\
Scope overlap & Shared or explicitly modelled geography, habitat, time, and inclusion rules & Domain totals reference different populations & Scope-specific interval table\\
Delimitation compatibility & Criteria and evidence thresholds are harmonized or their differences modelled & ``Species'' differs across components & Unit-specific richness ledger\\
Detection comparability & Observation processes and coverage are known, or adjustment is defensible & Apparent differences may reflect protocol rather than diversity & Method-stratified reporting\\
Uncertainty propagation & Mapping and overlap uncertainty enter the interval & Sum would be spuriously precise & No aggregate headline figure\\
Independent replication & Separate data/code workflow reproduces main decision & Integration rests on one pipeline or laboratory & Provisional status only\\
\bottomrule
\end{tabular}
\end{table*}

Table~\ref{tab:gates} converts a philosophical disagreement into a testable data-integration checklist. It has a practical consequence: a project may conclude that a named-taxon registry, a eukaryotic richness extrapolation, and a microbial ASV extrapolation are all credible for their own passports but must remain non-additive. This is an advance in reporting discipline, not a failure of synthesis.

\subsection{Validation targets}

Prediction is evaluated at two levels. For occurrence/detection predictions, let $\pi_r$ be a predicted probability for withheld replicated outcome $y_r$. The Brier score is
\begin{equation}
\mathrm{BS}=\frac{1}{R}\sum_{r=1}^{R}\left(y_r-\pi_r\right)^2,
\label{eq:brier}
\end{equation}
where $R$ is the number of held-out predictions. Equation~\eqref{eq:brier} tests calibration, not just discrimination. For richness intervals, a held-out stratum $h$ is evaluated by its coverage of an independently observed or cautiously extrapolated reference $S_h^{\mathrm{ref}}$:
\begin{equation}
I_h=\mathbb{I}\left\{S_h^{\mathrm{ref}}\in[L_h,U_h]\right\}, \qquad \widehat{\Gamma}=\frac{1}{H_{\mathrm{test}}}\sum_{h=1}^{H_{\mathrm{test}}}I_h.
\label{eq:interval}
\end{equation}
Equation~\eqref{eq:interval} estimates interval coverage across held-out strata, where $[L_h,U_h]$ is the planned interval and $\mathbb{I}$ an indicator. The reference itself can be uncertain; that uncertainty must be carried forward rather than treated as an oracle.

\section{Decision Schedule and Confirmatory Tests}

\subsection{From pilot to globally reusable modules}

The design is intentionally modular. A first implementation should not attempt to settle the entire Earth total. It begins with one bounded realm, a short list of declared units, and calibration sites where vouchers, expert identification, barcodes, and eDNA can be examined in the same time window. This pilot determines whether the observation protocol controls false detections, whether occupancy can be identified from the available replication, and whether coverage targets are attainable. Only a protocol that meets those criteria progresses to expanded strata.

The next module expands geographic and environmental strata while preserving the original passport fields. It adds a locked validation set before model selection and a registry of taxonomic concept changes. A third module permits cross-domain comparison only after the S4 gates are reviewed. Thus, scale is achieved by replication of a common interface, not by replacing locally interpretable estimates with a single opaque predictor. This schedule also directs investment toward the largest uncertainty component. For example, poor detection motivates field and laboratory replication, whereas unstable lineage delimitation motivates vouchering, reference sequencing, and expert revision rather than more samples of the same marker.

\begin{table*}[t]
\caption{Planned decision schedule. Advancement requires the stated evidence rather than merely additional data volume.}
\label{tab:schedule}
\centering
\footnotesize
\begin{tabular}{L{0.15\textwidth} L{0.23\textwidth} L{0.25\textwidth} L{0.25\textwidth}}
\toprule
\textbf{Module} & \textbf{Primary question} & \textbf{Confirmation evidence} & \textbf{Decision if criterion fails}\\
\midrule
Pilot calibration & Are detections and provisional units technically credible? & Controls pass; replicate detection is estimable; expert and molecular evidence are linked at calibration sites & Restrict to method development; no richness extrapolation\\
Within-domain expansion & Is a domain richness interval stable at target coverage? & Coverage diagnostics, sensitivity set, and held-out strata meet preregistered thresholds & Report only observed diversity or revise sample frame/model\\
Independent replication & Does a separate workflow reproduce the inference? & Distinct field/laboratory or analytic team reproduces key interval and S4 decision & Mark result provisional and diagnose source of discrepancy\\
Cross-domain integration & Is aggregation scientifically meaningful? & Unit crosswalk, overlap rule, compatible scope, and propagated uncertainty pass all gates & Publish a versioned vector, not a global sum\\
Update cycle & Does a new taxonomy or reference database change the claim? & Change log quantifies movement attributable to revision versus new observation & Release a new result version with comparability statement\\
\bottomrule
\end{tabular}
\end{table*}

Table~\ref{tab:schedule} gives the operational meaning of confirmatory progress. It avoids a common failure mode in global biodiversity work: adding heterogeneous records until an estimate appears precise while the data-generating process remains uncharacterized. A module can be successful even when it blocks advancement, because a documented failure identifies whether to redesign sampling, improve reference collections, increase technical controls, or narrow the stated estimand.

\subsection{Inference hierarchy}

The primary confirmatory endpoint is a coverage-adjusted interval for a preregistered domain and unit. Secondary endpoints are held-out detection calibration, interval coverage, and the uncertainty decomposition in Eq.~\eqref{eq:variance}. Comparisons among markers, clustering rules, and strata are treated as a family of tests. When formal null-hypothesis comparisons are appropriate, family-wise error is controlled across the preregistered family; when the objective is model assessment, calibration and predictive performance take precedence over a binary significance label. Exploratory discoveries remain explicitly exploratory until evaluated on independent material.

This hierarchy prevents the vocabulary of confirmation from being attached to a post hoc global extrapolation. It also makes negative results valuable. If a marker cannot detect a taxonomic group with usable sensitivity, or if a rare-tail model fails on held-out strata, the documented outcome improves the next sampling round without being transformed into a claim about biological absence. The final public report will distinguish detected units, supported lineage hypotheses, named taxa, and unresolved candidates in separate tables.

\section{Implementation and Data Governance}

\subsection{Reproducible workflow}

The proposed implementation has four immutable layers. The raw layer stores original reads, images, vouchers, collection events, permits, and control results. The curated layer stores taxonomic concept identifiers, synonym mappings, reference sequences, and versioned quality-control outputs. The analytical layer stores locked environments, code, model specifications, random seeds, and validation splits. The reporting layer stores estimand passports, intervals, sensitivity distributions, integration decisions, and machine-readable provenance. A change to a reference taxonomy or clustering algorithm generates a new versioned result, not an unannounced replacement.

All field and laboratory processes use predefined identifiers connecting physical samples to derived data. Negative controls are never discarded solely because they are inconvenient; their outcomes determine whether a batch passes the S1 evidence gate. Calibration sites with both expert inventory and molecular surveys are retained as long-term benchmarks. This arrangement is compatible with community-scale catalogues such as the Earth Microbiome Project while preserving local effort and method metadata~\cite{thompson2017}.

\begin{table}[t]
\caption{Failure-oriented reporting matrix.}
\label{tab:failure}
\centering
\footnotesize
\begin{tabular}{L{0.27\columnwidth} L{0.30\columnwidth} L{0.29\columnwidth}}
\toprule
\textbf{Observed failure} & \textbf{Interpretation} & \textbf{Required action}\\
\midrule
Controls detect target units & Possible contamination or carryover & Quarantine batch; diagnose before any richness update.\\
Coverage remains low & Observed count is incomplete for target frame & Expand replication or limit claim to observed diversity.\\
Threshold changes alter conclusion & Unit count is analysis-sensitive & Publish sensitivity; do not relabel as stable species total.\\
Held-out calibration fails & Model does not generalize to target conditions & Refit or downgrade inference to exploratory.\\
Crosswalk lacks evidence & Domains cannot be reconciled without assertion & Report estimates in parallel and identify bridge data needed.\\
Taxonomic revision changes count & Named-taxon estimate is version-dependent & Release a new passport and comparison log.\\
\bottomrule
\end{tabular}
\end{table}

\subsection{Ethics and governance}

Biodiversity measurement can create harm if it exposes locations of threatened taxa, overlooks Indigenous data governance, or treats samples as unowned inputs. The protocol therefore requires permits, access and benefit-sharing compliance, sensitive-location controls, community and Indigenous governance where applicable, and a separation between public aggregate outputs and restricted locality data. Species estimates also have distributional consequences: a claim that a region is ``data-poor'' may reflect unequal funding, language barriers, collection history, or access restrictions rather than low diversity. Sampling effort and missingness must be reported as scientific variables.

The design is compatible with the global biodiversity assessment need for better observation systems and curated knowledge, while rejecting the view that automation removes the need for taxonomic expertise~\cite{pimm2014,ipbes2019}. Automated discovery expands candidate space; curation, verification, and governance determine which statements can become durable biodiversity knowledge.

\section{Discussion}

The central scientific conclusion is conditional but strong: no unique total for ``all species on Earth'' is interpretable without a specification of what counts as a species and which part of life is within scope. This is not semantic evasion. It is the same measurement discipline that distinguishes a temperature from a heat content or a sequencing read from an organismal lineage. Once the estimand is defined, the problem becomes measurable and falsifiable: detection can be calibrated, delimitation can be challenged with new evidence, coverage can be increased, model extrapolations can be tested on held-out strata, and integration can be approved or declined by explicit gates.

\specscope{} changes the role of several familiar tools. Taxonomic catalogues become versioned evidence for named concepts rather than a proxy for all extant diversity. eDNA and metabarcoding become high-throughput discovery systems with controlled false-positive and false-negative pathways, not automatic species censuses. Coverage estimators become domain-specific completeness diagnostics, not black boxes that transform sparse samples into global certainty. Scaling laws become hypotheses with stated units and validation targets rather than universal conversion factors. These roles are mutually reinforcing precisely because they are not collapsed.

The approach also clarifies why a short public answer can be misleading. A number around several million may be a defensible statement about global eukaryotic taxonomic richness under a particular model~\cite{mora2011}. A number near one trillion may be a provocative but conditional statement about global microbial diversity under another model~\cite{locey2016}. Neither should be dismissed merely for being different; neither should be detached from its passport. The result of responsible synthesis may be an interval table or a vector of estimates rather than a single total.

There are limitations. The framework cannot make an inaccessible habitat sampled, create taxonomic capacity on demand, or resolve every lineage boundary. Occupancy models may be weakly identified with sparse replication; rare-tail estimators can remain sensitive where many units are singletons; and genomic, morphological, and ecological evidence can conflict. Integration gates may therefore reject a desired total for years. These are genuine scientific constraints, and making them visible helps allocate effort to calibration, reference building, vouchers, and under-sampled environments.

The next empirical study should preregister a restricted first target, such as a clearly bounded biome and a small set of organismal and molecular units. It should publish the S0 passport before field collection, run joint voucher-eDNA calibration sites, report coverage and sensitivity before extrapolating, and invite independent replication of both the data-processing workflow and the integration decision. A credible global programme can then grow through comparable modules rather than one opaque extrapolation.

\section{Conclusion}

Earth's species richness should not be represented as a contest between incompatible headline totals. The defensible answer is a family of scope-specific estimates, each accompanied by its species concept, counting unit, sampling frame, evidence protocol, coverage, estimator, uncertainty, and integration status. \specscope{} operationalizes this answer through five stages: declare scope, validate evidence, quantify coverage, calibrate estimation, and integrate only where the data support a common unit. This design does not claim a new total; it provides the pathway by which future totals can be tested, compared, and responsibly communicated.

\clearpage
\balance
\begin{thebibliography}{99}

\bibitem{may1988}
R. M. May, ``How many species are there on earth?'' \emph{Science}, vol. 241, no. 4872, pp. 1441--1449, 1988, doi: 10.1126/science.241.4872.1441.

\bibitem{mora2011}
C. Mora, D. P. Tittensor, S. Adl, A. G. B. Simpson, and B. Worm, ``How many species are there on Earth and in the ocean?'' \emph{PLoS Biology}, vol. 9, no. 8, e1001127, 2011, doi: 10.1371/journal.pbio.1001127.

\bibitem{locey2016}
K. J. Locey and J. T. Lennon, ``Scaling laws predict global microbial diversity,'' \emph{Proceedings of the National Academy of Sciences}, vol. 113, no. 21, pp. 5970--5975, 2016, doi: 10.1073/pnas.1521291113.

\bibitem{dequeiroz2007}
K. de Queiroz, ``Species concepts and species delimitation,'' \emph{Systematic Biology}, vol. 56, no. 6, pp. 879--886, 2007, doi: 10.1080/10635150701701083.

\bibitem{bickford2007}
D. Bickford \emph{et al.}, ``Cryptic species as a window on diversity and conservation,'' \emph{Trends in Ecology \& Evolution}, vol. 22, no. 3, pp. 148--155, 2007, doi: 10.1016/j.tree.2006.11.004.

\bibitem{leliaert2014}
F. Leliaert \emph{et al.}, ``DNA-based species delimitation in algae,'' \emph{European Journal of Phycology}, vol. 49, no. 2, pp. 179--196, 2014, doi: 10.1080/09670262.2014.904524.

\bibitem{fujisawa2013}
T. Fujisawa and T. G. Barraclough, ``Delimiting species using single-locus data and the generalized mixed Yule coalescent approach,'' \emph{Systematic Biology}, vol. 62, no. 5, pp. 707--724, 2013, doi: 10.1093/sysbio/syt033.

\bibitem{ratnasingham2013}
S. Ratnasingham and P. D. N. Hebert, ``A DNA-based registry for all animal species: The Barcode Index Number (BIN) system,'' \emph{PLoS ONE}, vol. 8, no. 7, e66213, 2013, doi: 10.1371/journal.pone.0066213.

\bibitem{froslev2017}
T. G. Frøslev \emph{et al.}, ``Algorithm for post-clustering curation of DNA amplicon data yields reliable biodiversity estimates,'' \emph{Nature Communications}, vol. 8, art. 1188, 2017, doi: 10.1038/s41467-017-01312-x.

\bibitem{deiner2017}
K. Deiner \emph{et al.}, ``Environmental DNA metabarcoding: Transforming how we survey animal and plant communities,'' \emph{Molecular Ecology}, vol. 26, no. 21, pp. 5872--5895, 2017, doi: 10.1111/mec.14350.

\bibitem{olds2016}
B. P. Olds \emph{et al.}, ``Estimating species richness using environmental DNA,'' \emph{Ecology and Evolution}, vol. 6, no. 12, pp. 4214--4226, 2016, doi: 10.1002/ece3.2186.

\bibitem{chao2012}
A. Chao and L. Jost, ``Coverage-based rarefaction and extrapolation: Standardizing samples by completeness rather than size,'' \emph{Ecology}, vol. 93, no. 12, pp. 2533--2547, 2012, doi: 10.1890/11-1952.1.

\bibitem{thompson2017}
L. R. Thompson \emph{et al.}, ``A communal catalogue reveals Earth's multiscale microbial diversity,'' \emph{Nature}, vol. 551, pp. 457--463, 2017, doi: 10.1038/nature24621.

\bibitem{pimm2014}
S. L. Pimm \emph{et al.}, ``The biodiversity of species and their rates of extinction, distribution, and protection,'' \emph{Science}, vol. 344, no. 6187, 1246752, 2014, doi: 10.1126/science.1246752.

\bibitem{ipbes2019}
IPBES, \emph{Global Assessment Report on Biodiversity and Ecosystem Services}. Bonn, Germany: IPBES Secretariat, 2019, doi: 10.5281/zenodo.3831673.

\end{thebibliography}

\end{document}
